\documentclass[12pt]{article}
\usepackage{amsmath,amssymb}

\begin{document}
\section*{{2024 AMC 12A Problems/Problem 18}}
Source: AoPS Wiki

\subsection*{Problem}
On top of a rectangular card with sides of length $1$ and $2+\sqrt{3}$ , an identical card is placed so that two of their diagonals line up, as shown ( $\overline{AC}$ , in this case). [asy] defaultpen(fontsize(12)+0.85); size(150); real h=2.25; pair C=origin,B=(0,h),A=(1,h),D=(1,0),Dp=reflect(A,C)*D,Bp=reflect(A,C)*B; pair L=extension(A,Dp,B,C),R=extension(Bp,C,A,D); draw(L--B--A--Dp--C--Bp--A); draw(C--D--R); draw(L--C^^R--A,dashed+0.6); draw(A--C,black+0.6); dot("$C$",C,2*dir(C-R)); dot("$A$",A,1.5*dir(A-L)); dot("$B$",B,dir(B-R)); [/asy] Continue the process, adding a third card to the second, and so on, lining up successive diagonals after rotating clockwise. In total, how many cards must be used until a vertex of a new card lands exactly on the vertex labeled $B$ in the figure? $\textbf{(A) }6\qquad\textbf{(B) }8\qquad\textbf{(C) }10\qquad\textbf{(D) }12\qquad\textbf{(E) }\text{No new vertex will land on }B.$

\end{document}
